% 2013 年北京卷理科第 4 题：程序框图。
% 来源 PDF：普通高考/2013/2013北京理.pdf，第 1 页（印面 787）左栏。
% 流程：开始 → i=0,S=1 → S=(S^2+1)/(2S+1) → i=i+1 → i≥2?
%       是：输出 S → 结束；否：回到 S 赋值框。
\begin{tikzpicture}[font=\footnotesize, >=Stealth, line width=0.45pt,
  node distance=4.5mm,
  term/.style={draw, rounded corners=6pt, minimum height=1.8em, inner xsep=9pt},
  io/.style={draw, trapezium, trapezium left angle=70, trapezium right angle=110,
             minimum height=1.8em, inner xsep=6pt},
  box/.style={draw, minimum height=1.9em, inner xsep=6pt},
  dec/.style={draw, diamond, aspect=2.2, inner sep=1pt}]
  \node[term] (start) {开始};
  \node[box, below=4mm of start] (init) {$i=0$，$S=1$};
  \node[box, below=4mm of init] (assign) {$S=\dfrac{S^2+1}{2S+1}$};
  \node[box, below=4mm of assign] (inc) {$i=i+1$};
  \node[dec, below=5mm of inc] (cond) {$i\geqslant 2$};
  \node[io, right=16mm of cond] (output) {输出 $S$};
  \node[term, below=5mm of output] (end) {结束};
  \draw[->] (start) -- (init);
  \draw[->] (init) -- (assign);
  \draw[->] (assign) -- (inc);
  \draw[->] (inc) -- (cond);
  \draw[->] (cond) -- node[above] {是} (output);
  \draw[->] (output) -- (end);
  % 循环支路：从判断框左侧引出，向上回到 S 赋值框左侧
  \draw[->] (cond.west) -- node[above] {否} ++(-2.1,0) |- (assign.west);
\end{tikzpicture}
